\par\clearpage\SourceLeaf{284}\ArithmeticHead{199}{R}{Xa. DIE JUBILÄUMSSCHRIFT.}
\begin{center}
\textbf{X. Zunehmendes Alter. Doktorjubiläum. Letzte Lebensjahre.}\par
\textls[100]{a) Jubiläumsschrift.}\par
\textls[100]{Plan einer Abhandlung über Reihenkonvergenz.}
\end{center}
\noindent\hspace*{\parindent}\textsc{Sartorius} berichtet\footnote{A. a. O. Fussnote \textsuperscript{13)}, S. 78.}, \textsc{Gauss} habe einmal gesagt, dass ihm in seiner
Jugend die Gedanken in solcher Fülle ununterbrochen zugeströmt seien, dass
er ihrer kaum hätte Herr werden, und nur einen Teil derselben hätte
aufzeichnen können. Überblickt man aber seine Produktion auch in der
späteren Zeit, also etwa von 1840 ab, so hat man den Eindruck, dass dieselbe
Ideenfülle noch bis ins hohe Alter vorgehalten hat, und nur infolge der
Abnahme der körperlichen Kräfte nicht mehr voll zur Gestaltung kommen
konnte. Wir nennen unter den auch durch ihre Vielseitigkeit erstaunlichen
Arbeiten dieser Zeit die \textit{Dioptrischen Untersuchungen} (1840, Werke V, S. 243,
\textit{Anzeige} S. 305), die beiden Abhandlungen \textit{über höhere Geodäsie} (I, 1843,
Werke IV, S. 259, \textit{Anzeige} S. 347; II, 1846, ebenda, S. 301, \textit{Anzeige} S. 352)
und die grosse, in das Jahr 1845 fallende Denkschrift über die Neuorganisation
der Professoren-Witwenkasse (abgedruckt Werke IV, S. 119). Aus dem
Jahre 1847 sind uns drei in das Gebiet der Analysis gehörende Aufzeichnungen
erhalten, die sich auf den letzten Seiten des Handbuchs 19, Be finden,
es sind dies
\begin{list}{}{\setlength{\leftmargin}{2em}\setlength{\labelwidth}{1.5em}\setlength{\labelsep}{.5em}\setlength{\itemsep}{0pt}\setlength{\parsep}{0pt}\setlength{\topsep}{0pt}}
\item[1)] S. 273—276 des Handbuchs (abgedruckt Werke VIII, S. 80—83) »\textsc{Lagranges} Lehrsatz auf möglichst lichtvolle Art abgeleitet«, datiert vom
13. Mai 1847, vergl. oben S. 25.
\item[2)] S. 277 des Handbuchs (abgedruckt Werke X, 1, S. 398, 399) die Darstellung
einiger diskontinuierlicher Funktionen durch trigonometrische
Reihen,
\item[3)] S. 289—293 des Handbuchs (abgedruckt Werke X, 1, S. 355—359) einige
Ansätze zur Theorie der Reihen \(F(\alpha,\beta,\gamma,x)\), insbesondere derjenigen, die
(vergl. Art. 6 der \textit{Disquisitiones} von 1812) als Entwicklungskoeffizienten
von \((1+f^2-2f\cos\varphi)^{-\theta}\) nach den Kosinus der Vielfachen von \(\varphi\) auftreten.
\textsc{Gauss} verfolgt hier anscheinend den Zweck, das Verhalten dieser
Koeffizienten für weit entfernte Glieder zu bestimmen, ähnlich wie er
\end{list}
\clearpage\SourceLeaf{285}\ArithmeticHead{200}{L}{L. SCHLESINGER, ÜBER GAUSS' FUNKTIONENTHEORETISCHE ARBEITEN.}
\begingroup\leftskip=2em\noindent
es in früher Jugendzeit für die Entwicklungskoeffizienten der Mittelpunktsgleichung
getan hat (siehe Abschn. Ve), oben S. 115).\par\endgroup
\noindent Mit dieser letztgenannten Aufzeichnung schliesst das im Mai 1809 begonnene
Handbuch 19, Be ab; \textsc{Gauss} hat danach kein neues Handbuch mehr angefangen,
sondern seine mathematischen Notizen und Entwürfe auf lose Blätter
geschrieben.
\par
Das Jahr 1849 bringt am 16. Juli \textsc{Gauss'} goldenes Doktorjubiläum\footnote{Die promovierende Universität Helmstedt war 1809 durch die westfälische Regierung \textsc{Jérômes} aufgehoben worden.},
das der Jubilar selbst dadurch feiert, dass er in der festlich geschmückten
Aula der Gesellschaft der Wissenschaften die Abhandlung \textit{Beiträge zur Theorie
der Gleichungen} (Werke III, S. 71, \textit{Anzeige} S. 113) vorlegte. Der erste Teil
dieser Arbeit bezweckt, den vor fünfzig Jahren gegebenen Beweis für die
Wurzelexistenz »mit erheblichen Zusätzen« und »in einer abgeänderten und
\mbox{.\,.\,.} eine vergrösserte Klarheit darbietenden Gestalt zu wiederholen«, während
der zweite Teil Methoden angibt, um für trinomische Gleichungen »nicht
bloss die reellen, sondern auch die imaginären Wurzeln \mbox{.\,.\,.} mit Leichtigkeit
zu bestimmen«.
\par
Die »erheblichen Zusätze«, die der erste Teil der Jubiläumsschrift zu dem
Beweise der Inauguraldissertation bringt, bestehen darin, dass einmal die Koeffizienten
der gegebenen ganzen Funktion \(n\)-ten Grades \(X=T+Ui\) beliebig
komplex sein können, und dass ferner die Existenz der \(n\) Wurzeln auf einen
Schlag bewiesen wird, was (vergl. oben S. 61) ja schon im Art. 23 der \textit{Dissertation}
als möglich bezeichnet worden war. Die Abänderung und Vereinfachung
der Beweisführung erzielt \textsc{Gauss} (vergl. die Fussnote Werke III, S. 82)
dadurch, dass er durch Betrachtung der Ebenenteile, in denen \(T\) positiv ist,
zeigt, dass die Kurve \(T=0\) in den Kreis \(K\), ausserhalb dessen sicher keine
Wurzel von \(X=0\) liegt, in \(n\) getrennten Zügen eintritt, von denen jeder
das Kreisinnere stetig durchsetzt und wieder verlässt. An der Eintritts- und
Austrittsstelle eines jeden Zuges hat \(U\), entgegengesetzte Vorzeichen, so dass
auf jedem eine Nullstelle von \(U\), d. h. also ein Punkt liegen muss, in dem \(T\)
und \(U\) gleichzeitig verschwinden. Im Art. 10 wird auch der Fall mehrfacher
Wurzeln erledigt. Die Betrachtung der Flächenräume, wo \(T>0\) ist, hat
\textsc{Gauss} seiner neuen Darstellung zugrunde gelegt, weil ihm die Natur der Ge-\par
\clearpage\SourceLeaf{286}\ArithmeticHead{201}{R}{Xa. DIE JUBILÄUMSSCHRIFT.}
\noindent biete, in denen eine Potentialfunktion konstantes Vorzeichen bewahrt, vertraut
geworden war. Er macht zwar in der Abhandlung selbst von dieser Eigenschaft
der Funktion \(T\) keinen direkten Gebrauch; wenn er aber im Art. 8
sagt, es »würde nicht schwer sein zu beweisen, dass bei der besonderen Beschaffenheit
der Funktion \(T\)« die Fläche, wo \(T>0\) ist, keine »nicht zu ihr
gehörige Insel« einschliessen kann, so hat er dabei ohne Zweifel die Beschaffenheit
von \(T\), ein logarithmisches Potential zu sein, im Sinne.
\par
Im zweiten Teile wird zunächst eine bis ins Einzelne gehende Vorschrift
entwickelt, wie für eine trinomische Gleichung
\[
x^{m+n}+fx^m+g=0
\]
mit reellen \(f,g\) die reellen Wurzeln mit Benutzung der \textsc{Gaussschen} Logarithmen
(vergl. oben Abschn. VIIe), S. 155) berechnet werden können. \textsc{Gauss}
hebt dabei (Werke III, S. 86) hervor, dass er dieses Verfahren schon 1840
für den Fall quadratischer Gleichungen (Werke III, S. 255) und 1843 für die
kubische Gleichung, die »bei der parabolischen Bewegung zur Bestimmung
der wahren Anomalie dient« (Werke VI, S. 191—193) angegeben, und an der
letztgenannten Stelle auf seine Anwendbarkeit bei beliebigen trinomischen
Gleichungen hingewiesen habe\footnote{Eine Erweiterung des \textsc{Gaussschen} Verfahrens gibt \textsc{Gundelfinger} in den \textit{Tafeln zur Berechnung der reellen Wurzeln sämtlicher Gleichungen etc.}, Leipzig 1897.}. Die Berechnung der komplexen Wurzeln
kann für eine Gleichung der Form
\[
x^{m+n}+fe^{\varphi i}x^m+ge^{\gamma i}=0
\]
von dem zweidimensionalen Gebiet auf das eindimensionale reduziert werden,
indem man \(x=re^{\rho i}\) setzt und \(\lambda=g^nf^{-n-m}\) einführt, wodurch sich zur Bestimmung
von \(\rho\) die Gleichung
\[
\lambda=(-1)^{m+n}\frac{\sin^m(m\rho+\varphi-\gamma)\sin^n(n\rho-\varphi)}{\sin^{m+n}((m+n)\rho-\gamma)}
\]
ergibt. Hat man hieraus \(\rho\) gefunden, so bestimmt sich \(r\) z. B. aus der Gleichung
\[
r^n=-\frac{f\sin(m\rho+\varphi-\gamma)}{\sin((m+n)\rho-\gamma)}.
\]
Hätte \textsc{Gauss} seine in dem Briefe an \textsc{Moebius} vom 13. August 1849 (Werke
X, 1, S. 109, siehe oben S. 190) ausgesprochene Absicht, in einem dritten
Teile der Jubiläumsschrift die Konfiguration der Kurven \(U=\mathrm{const},\ T=\mathrm{const}\)
zu untersuchen, durchgeführt, so würde er wohl auch das für die trinomischen\par
\HistoricalFoot{X 2 Abh. 2.}{26}
\clearpage\SourceLeaf{287}\ArithmeticHead{202}{L}{L. SCHLESINGER, ÜBER GAUSS' FUNKTIONENTHEORETISCHE ARBEITEN.}
\noindent Gleichungen entwickelte Verfahren dem oben skizzierten allgemeinen Gedankengang
eingeordnet haben.
\par
\textsc{Sartorius} berichtet (a. a. O. \textsuperscript{13)}, S. 69), \textsc{Gauss} habe nach dem Jubiläum öfter
erklärt, er möchte sich in seinen Arbeiten nicht treiben lassen, seine Arbeitszeit
sei im Vergleich zu früheren Jahren merklich kürzer geworden. Nach
jener Zeit sei \textsc{Gauss} hauptsächlich mit drei Gegenständen beschäftigt gewesen,
mit der Theorie der Konvergenz der Reihen, mit einer zweiten Revision und
Bilanz der Professoren-Witwenkasse und mit mechanischen Problemen über die
Erdrotation, im Anschluss an den \textsc{Foucaultschen} Pendelversuch. Über diese
letztgenannte Untersuchung vergleiche man das Werke XI, 1, S. 33—45 abgedruckte
Briefmaterial. Für die geplante Abhandlung zur Konvergenz der
Reihen sind im Nachlass (vergl. oben S. 142) fünf verschiedene Fassungen
ihres Anfangs vorhanden, über die Werke X, 1, S. 437 berichtet ist, und von
denen ebenda, S. 400—419 die beiden ersten Artikel des vierten Entwurfs
und der fünfte unverkürzt wiedergegeben sind. Davon, was \textsc{Gauss} beabsichtigt,
geben diese Entwürfe nur ein sehr unvollkommenes Bild, etwas mehr
erfahren wir aus einer Briefstelle an \textsc{Schumacher} vom 5. Februar 1850 (Werke
X, 1, S. 431). \textsc{Gauss} schreibt, er sei nicht abgeneigt, eine ihm zuteil werdende
Musse dazu zu verwenden, um seine Methode, »den Grad der Konvergenz der
nach Kosinus und Sinus der Vielfachen eines Winkels fortschreitenden, eine
beliebige periodische Funktion ausdrückenden Reihe zu bestimmen, \mbox{.\,.\,.\,.} in
einer ihm selbst genügenden Gestalt auszuführen«, und nennt als hierhergehörig,
die Untersuchung über die Mittelpunktsgleichung, über die oben
(Abschn. Ve), S. 115ff.) berichtet worden ist. Aus diesem Hinweis geht schon
hervor, was \textsc{Gauss} hier unter Konvergenz versteht; der vierte Entwurf, betitelt
\textit{Konvergenz der Reihen, in welche die periodischen Funktionen einer
veränderlichen Grösse entwickelt werden} gibt in seinem ersten Artikel, der
die Überschrift trägt »Konvergenz der unendlichen Reihen im Allgemeinen«,
ausdrücklich die folgende Definition: »Ich werde unter Konvergenz, einer
unendlichen Reihe schlechthin beigelegt, nichts anderes verstehen, als die beim
unendlichen Fortschreiten der Reihe eintretende unendliche Annäherung ihrer
Glieder an die Null«\footnote{Vergl. demgegenüber den Gebrauch des Wortes Konvergenz in dem uns geläufigen Sinne in dem Briefe an \textsc{Schumacher} vom 1. September 1850, Werke X, 1, S. 435.}. Im Art. 1 des fünften Entwurfs (Werke X, 1, S. 407)\par
\clearpage\SourceLeaf{288}\ArithmeticHead{203}{R}{Xa. ABHANDLUNG ÜBER REIHEN.}
\noindent heisst es dann: »Die Überschrift bezeichnet zwar den Hauptgegenstand dieser
Denkschrift, der jedoch nur einen Teil ihres Inhalts ausmacht. Die Lösung
der Aufgabe gehört recht eigentlich in das Gebiet der Lehre von den komplexen
Grössen \mbox{.\,.\,.\,.}; allein die Wege in diesem Gebiet sind noch nicht
überall gebahnt. \mbox{.\,.\,.} Es werden daher mit dem Hauptgegenstande einige
Nebenuntersuchungen \mbox{.\,.\,.} verbunden werden müssen«. Die Artikel 2—5 geben
dann für durch »Züge« begrenzte »Schichten« in der Ebene der komplexen
Variabelen die Grundlagen der Topologie, wobei unter Bezugnahme auf den
Art. 5 der Jubiläumsschrift gesagt wird, dass diese geometrische Deutung »nur
bezweckt, die Bewegung in dem \mbox{.\,.\,.} Felde der abstrakten komplexen Grössen
zu erleichtern und eine Sprache für dieselbe zu vermitteln«. Eine Schicht
kann durch einen oder auch durch mehrere geschlossene Züge begrenzt werden,
und bestimmt allemal eine gewisse »Folgeordnung« in diesen Zügen, »ähnlich,
als wenn die sämtlichen Uferlinien eines eine oder mehrere Inseln einschliessenden
Sees in einem solchen Sinne durchlaufen werden, dass man \mbox{.\,.\,.}
den See \mbox{.\,.\,.} entweder allemal rechts oder allemal links hat«. Im Art. 5
wird dann gezeigt, »dass eine Schicht \(\Sigma\), die durch zwei oder mehrere geschlossene
Züge begrenzt wird, sich in ebensoviele partielle Schichten zerlegen
lässt, die jede nur von einem Zuge begrenzt werden«; wenn es sich dann um
Anwendungen handelt, »in welchen zwei an sich gleiche, aber in entgegengesetzten
Folgeordnungen vorkommende Züge wie einander destruierend betrachtet
werden können, kann man die Totalität der Grenzzüge von \(\Sigma\) \mbox{.\,.\,.}
gleich setzen \mbox{.\,.\,.} der Totalität der einfachen Grenzzüge sämtlicher aus der
\mbox{.\,.\,.} Zerlegung hervorgehender Schichten«\footnote{Vergl. in bezug auf \textsc{Gauss'} topologische Untersuchungen \textsc{Stäckel}, Abhandl. 4 dieses Bandes S. 49—56.}. Welche Anwendungen \textsc{Gauss}
hier im Auge hatte, ist nicht schwer zu erraten, geht aber mit Sicherheit
aus dem Inhalt der nachfolgenden Art. 6—8 hervor. Im Art. 6 gibt \textsc{Gauss}
die Definition eines Integrals \(\displaystyle\int\limits_{t_0}^T f(t)\,dt\) für eine zwischen \(t_0\) und \(T\) eindeutige,
endliche und stetige Funktion \(f(t)\) als Grenzwert einer Summe und spricht
sich dann (in den Artikeln 7, 8) über die Bedeutung eines solchen Integrals
weiter wie folgt aus:
\par
»Ist hingegen \(f(t)\) eine vielwertige Funktion, so ist erforderlich, dass nach\par
\HistoricalFoot{}{26\textsuperscript{*}}
\clearpage\SourceLeaf{289}\ArithmeticHead{204}{L}{L. SCHLESINGER, ÜBER GAUSS' FUNKTIONENTHEORETISCHE ARBEITEN.}
\noindent einem Prinzipe bestimmt sei, welcher aus den verschiedenen Werten der Funktion
für jeden Wert von \(t\) zwischen den Grenzen \(t_0\) und \(T\) gelten soll. Das
Prinzip, dass die Werte der Funktion nach der Stetigkeit zusammenhängen
sollen, wird, sobald \(f(t)\) für einen der Werte \(t\) gegeben ist, zur Bestimmung des
Zuges ausreichen, so lange man nicht auf einen Wert von \(t\) kommt, für
welchen der betreffende Wert von \(f(t)\) mit einem andern zusammenfällt (also
zwei Züge sich kreuzen), in welchem Falle über das fernere Fortschreiten eine
anderweitige Bestimmung getroffen werden muss. Übrigens verstehe ich das
Fortschreiten nach der Stetigkeit nur so, dass unendlich kleinen Veränderungen
von \(t\) unendlich kleine Änderungen von \(f(t)\) entsprechen, und sehe einen Übergang
von reellen Werten zu imaginären \textls[120]{nicht} als eine Unterbrechung der
Stetigkeit an.«
\par
Über den Fall, wo \(f(t)\) als eindeutige Funktion in einem zwischen \(t_0\) und
\(T\) gelegenen Werte \(t'\) unendlich wird, heisst es weiter:
\par
»Wenn, indem \(f(t')\) unendlich wird, \(\displaystyle\int\limits_{t_0}^{t'-\omega}f(t)\,dt\) und \(\displaystyle\int\limits_{t'+\omega}^T f(t)\,dt\), bei unendlich
abnehmendem \(\omega\), nicht beide endliche Grenzwerte haben, so bildet \mbox{.\,.\,.} der Wert
\(t=t'\) eine unübersteigliche Scheidewand, so lange man (wie man bisher immer
getan hat) sich dabei ausschliesslich auf reelle Werte von \(t\) beschränkt. Es
findet sich hier also eine Lücke in der Integralrechnung. \mbox{.\,.\,.} Die gründliche
Abhülfe dieses Mangels ist nur dadurch zu gewinnen, dass man den imaginären
Grössen völlig gleiche Rechte einräumt, also die Analysis gleichmässig
über das ganze Gebiet der komplexen Grössen erstreckt. In diesem Gebiete
kann man von einem Werte der veränderlichen Grösse \(t\) zu einem andern
auf unendlich vielen verschiedenen Wegen nach der Stetigkeit gelangen, und
so bei der Integration \(\int f(t)\,dt\) solche Werte von \(t\), für welche \(f(t)\) unendlich
gross wird, umgehen. Eine vollständige Abhandlung der aus diesem Gesichtspunkte
aufgefassten Theorie des Integrierens würde einen viel grösseren Raum
erfordern, als hier zulässig ist. Es werden jedoch einige für unsere Hauptuntersuchung
notwendigen Sätze hier entwickelt werden müssen.«
\par
Damit bricht die Handschrift ab.
\par
\textsc{Gauss} hatte also die Absicht, die in den Art. 2—5 entwickelte Topologie
der schlichten ebenen Bereiche auf die Theorie der Integration im komplexen
Gebiete anzuwenden, wobei jedoch die Beschränkung auf eindeutige\par
\clearpage\SourceLeaf{290}\ArithmeticHead{205}{R}{Xa. ABHANDLUNG ÜBER REIHEN.}
\noindent Funktionen, bezw. eindeutige Zweige mehrdeutiger Funktionen wesentlich erscheint.
\par
Das, was \textsc{Gauss} über die Integrale \textls[120]{mehrdeutiger} Funktionen sagt, bestätigt
nämlich (vergl. oben S. 152ff.), dass er in das Wesen dieser Integrale
nicht völlig eingedrungen ist. Denn dadurch, dass er nach dem Einrücken
in einen Verzweigungspunkt ȟber das fernere Fortschreiten eine anderweitige
Bestimmung« für erforderlich hält, beschränkt er sich auf die Betrachtung
eines eindeutigen Zweiges der Integralfunktion, statt — wie es z. B. \textsc{Jacobi}
1835\footnote{Vergl. Bibliotheca Mathematica (III), 9, 1909, S. 211—226 (\textsc{Gundelfinger}); 11, 1911, S. 138—152 (\textsc{Schlesinger}).} getan hat —, gerade aus dem Fehlen einer solchen Bestimmung die
Abhängigkeit des Integralwerts vom Integrationswege zu folgern.
\par
In bezug auf den »Hauptgegenstand der Denkschrift« haben wir nur noch
eine (als Art. 9, Werke X, 1, S. 416—419 abgedruckte) Aufzeichnung auf
einem einzelnen Blatte, und die offenbar gleichfalls hierhergehörigen, oben
(S. 199) unter 1)—3) genannten letzten Eintragungen in dem Handbuch 19, Be.
Ein deutliches Bild von dem, was \textsc{Gauss} beabsichtigt hat, lässt sich aber auch
aus diesen Notizen nicht gewinnen. Weitere Aufzeichnungen aus dem Gebiete
der Analysis sind uns nicht erhalten.
\par
An kleineren Bemerkungen stammen aus 1850 noch die im Handbuch
von \textsc{Rümcker} veröffentlichte geometrische Aufgabe (Werke IV, S. 407—412),
die in Briefen an \textsc{Schumacher} gegebene Lösung des Achtköniginnenproblems
(Werke XII, S. 19—27) und der Brief an den Rechenkünstler \textsc{Dase} (ebenda,
S. 44—47), aus 1852 die Lösung einer Aufgabe aus der Illustrierten Zeitung
(ebenda, S. 10—11).
\par
Wenn wir zusammenfassen, wie weit \textsc{Gauss} in das Gebiet der allgemeinen
Theorie der Funktionen komplexer Variabeln vorgedrungen war, so können
wir sagen, dass er die Mehrdeutigkeit einer solchen Funktion durch ihre Wertänderung
bei einem Umlauf (reditus) der Variabeln um einen singulären Punkt
erklärt und die Integration in einem Gebiete, wo die Funktion eindeutig ist,
beherrscht hat. Er hat ferner die Bedeutung der Topologie schlichter ebener
Bereiche für die Theorie dieser Integrale gekannt, und den Zusammenhang
entdeckt, der zwischen den Funktionen einer komplexen Variabeln, der konformen
Abbildung und der Theorie des logarithmischen Potentials besteht.\par
\clearpage\SourceLeaf{291}\ArithmeticHead{206}{L}{L. SCHLESINGER, ÜBER GAUSS' FUNKTIONENTHEORETISCHE ARBEITEN.}
In bezug auf die Einsicht in diese Zusammenhänge, sowie auch in bezug
auf die Methoden der geometrischen Funktionentheorie, die \textsc{Gauss} schon in
ganz früher Zeit gekannt und mit Erfolg angewandt hat, ist er \textsc{Cauchy} und
dessen Schülern dauernd überlegen geblieben. Was er durch die für alle
Gebiete der Grössenlehre geforderte und in seinen Untersuchungen zur Geltung
gebrachte Strenge in den Grundlagen und in der Beweisführung — von
ihm als Rigor antiquus bezeichnet — als Reformator der ganzen Mathematik
geleistet hat, kommt sehr prägnant in der Adresse zum Ausdruck, die ihm
die Berliner Akademie zum Doktorjubiläum gewidmet hat\footnote{Bericht über die Verhandlungen der Kgl. Preuss. Akad. d. W. zu Berlin 1849, S. 207—209. Die Urschrift befindet sich im Gaussarchiv, sie ist von sämtlichen Mitgliedern unterzeichnet, darunter von den Mathematikern \textsc{Jacobi}, \textsc{Lejeune Dirichlet} und \textsc{Steiner}.}. »In allen diesen
Schriften«, heisst es daselbst, »gesellt sich zu dem hohen Werte des sachlichen
Inhalts die vollendete Form der Darstellung, und eine solche Strenge der
Beweisführung, dass Sie die Behandlung mathematischer Gegenstände, welche
bei Ihren Vorgängern nicht selten in Mechanismus auszuarten schien, wieder
auf die antike Strenge zurückgebracht haben, und dass in Ihren Werken
nirgends Resultate auftreten, die nur mit gewissen Beschränkungen gelten, ohne
dass diese mit Notwendigkeit aus der Betrachtungsweise hervorgehen, wodurch
Sie die Resultate begründen«. Die Grenzen, die seinen Einsichten gesetzt
waren, hat \textsc{Gauss} selbst am schärfsten gesehen und am deutlichsten gefühlt;
darum liess er in seinen Handbüchern Leistungen verborgen, zu deren Verständnis
die nach ihm Gekommenen nur allmählich, oft erst nach Jahrzehnten
herangereift sind.
